%=============================================================================!
% 15.1  THE ERROR OF AN AVERAGE                                               !
%=============================================================================!
\section{The error of an average}
\label{sec:cv-mean}

A run reports the average of a quantity over its samples, and that average
misses the ensemble average by an amount that changes from run to run.
This section finds the size of the miss when the samples are independent,
on rolls of a die first, and says what an error bar built from it
promises.

\subsection*{The mean of independent rolls}

A fair die shows each face with probability $\frac16$, so one roll $x$ has
the mean $\ens{x} = 3.5$ and the variance $\sigma_x^2 = 35/12$
(\cref{ex:sm-die}). Rolling it $n$ times and averaging gives
\begin{equation*}
   \bar x = \frac1n\sum_{i=1}^nx_i ,
\end{equation*}
the overbar marking an average over samples, as it marks an average over
time (\cref{sec:sm-ergodic}). Each roll has the mean $\ens{x}$, and the
expectation of a sum is the sum of the expectations, so $\ens{\bar x} =
n\ens{x}/n = \ens{x}$: many such averages centre on the right value.
Their spread follows from the rule of \cref{sec:sm-probability} that the
variances of independent quantities add, and from one more. The rolls are
independent, so the variances of the $n$ of them add,
$\operatorname{var}(\sum_ix_i) = n\sigma_x^2$. A constant factor $a$
multiplies a variance by $a^2$, since
\begin{equation*}
   \operatorname{var}(ax) = \ens{a^2x^2} - \ens{ax}^2
   = a^2\bigl(\ens{x^2} - \ens{x}^2\bigr) = a^2\operatorname{var}(x) ,
\end{equation*}
and with $a = 1/n$,
\begin{equation}
   \operatorname{var}(\bar x) = \frac{n\sigma_x^2}{n^2}
   = \frac{\sigma_x^2}{n} , \qquad
   \sigma_{\bar x} = \frac{\sigma_x}{\sqrt n} .
   \label{eq:cv-mean}
\end{equation}
The spread of the mean falls as the square root of the number of samples:
four times as many rolls halve it, and a hundred times as many cut it to a
tenth. \Cref{fig:cv-dice}(b) rolls the die by computer, $n$ rolls at a
time, 200\,000 times over: the spread of the mean is 0.8558 for 4 rolls
against $\sigma_x/\sqrt4 = 0.8539$, and 0.1067 for 256 rolls against
0.1067.

\Cref{fig:cv-dice}(a) shows the shape as well as the spread. One roll has
six equally likely values, but the mean of 4 rolls is already close to
the Gaussian of \cref{eq:sm-gauss} with mean 3.5 and variance $\sigma_x^2/4$,
and the mean of 64 is indistinguishable from its Gaussian. The mean of
many independent samples approaches a Gaussian density whatever the
density of one sample, provided that density has a finite variance, a
result known as the central limit theorem, which the figure illustrates
without proving\cite{Grossfield2018}.

\begin{figure}[htb]
\centering
\includegraphics{ch15_convergence/dice}
\caption{The error of an average of independent rolls of a die, each set
of $n$ rolls repeated 200\,000 times. (a) The density of the mean of 4
rolls (light) and of 64 (dark), against the Gaussians of mean 3.5 and
variance $\sigma_x^2/n$ (dashed). (b) The spread of the mean against $n$,
against $\sigma_x/\sqrt n$ (dashed). (c) The fraction of the intervals
$\bar x \pm 1.96\,s/\sqrt n$ that hold 3.5 (teal), and of those with
Student's factor in place of 1.96 (ochre), against 0.95 (dotted).}
\label{fig:cv-dice}
\end{figure}

\subsection*{The standard error}

In a simulation $\sigma_x$ is not known in advance and must be estimated
from the same samples. The spread of the samples about their own mean
falls short of $\sigma_x$, because $\bar x$ sits closer to the samples than
$\ens{x}$ does. Writing each deviation as $x_i - \bar x = (x_i - \ens{x})
- (\bar x - \ens{x})$ and expanding the square,
\begin{align*}
   \sum_i(x_i - \bar x)^2
   &= \sum_i(x_i - \ens{x})^2 - 2(\bar x - \ens{x})\sum_i(x_i - \ens{x})
   + n(\bar x - \ens{x})^2\\
   &= \sum_i(x_i - \ens{x})^2 - n(\bar x - \ens{x})^2 ,
\end{align*}
since $\sum_i(x_i - \ens{x}) = n(\bar x - \ens{x})$. Averaging, the first sum gives
$n\sigma_x^2$ and the second term $n\operatorname{var}(\bar x) =
\sigma_x^2$, so the sum of squares averages to $(n - 1)\sigma_x^2$.
Dividing by $n - 1$ rather than $n$ therefore gives an estimate whose
average is $\sigma_x^2$,
\begin{equation*}
   s^2 = \frac{1}{n - 1}\sum_i(x_i - \bar x)^2 ,
\end{equation*}
and the estimate of $\sigma_{\bar x}$ is $s/\sqrt n$, the standard error of
the mean. The book writes a result as $\bar x \pm s/\sqrt n$, one standard
error, unless a factor is stated.

\subsection*{What an error bar promises}

When the mean is Gaussian, it lies within one $\sigma_{\bar x}$ of
$\ens{x}$ in 68.27\% of runs and within two in 95.45\%
(\cref{sec:sm-probability}); the factor that encloses 95\%, found
numerically from $\operatorname{erf}(1.960/\sqrt2) = 0.95$, is 1.960. An interval $\bar x \pm
1.96\,s/\sqrt n$ should therefore hold the true mean in 95\% of runs, and
\cref{fig:cv-dice}(c) counts how often it does: 0.949 of the time for 256
rolls and 0.905 for 8; for 3 rolls, not drawn, it is 0.777. With few
samples the estimate $s$ is itself uncertain, and a run whose $s$ happens
to be small claims an interval too narrow. For samples with a Gaussian
density the remedy is exact: Student showed that $(\bar x - \ens{x})/(s/
\sqrt n)$ then has a known density that depends on $n$ only through $n -
1$, its number of degrees of freedom. The $n$ deviations $x_i - \bar x$
from which $s$ is built add to zero, which leaves $n - 1$ of them free, as
holding the total momentum at zero left $3N - 3$ velocities free
(\cref{sec:sm-equipartition}). The density's 95\% factor is 4.303 for
three samples, 2.365 for 8 and 1.998 for 64, and approaches 1.960 as $n$
grows\cite{Student1908}. Three Gaussian samples with the factor 1.96
hold their mean in 0.811 of trials, as Student's density predicts; with
4.303 in 0.950. The die is far from Gaussian, and Student's factor brings
its three rolls only to 0.917. The standard errors of three runs that
Chapter~13 quoted are of this kind. For the difference of two means of
three runs each, Student's density has, when the two spreads are alike,
$2 + 2 = 4$ degrees of freedom and the 95\% factor 2.776, so the
differences of up to 2.5 combined errors that \cref{sec:th-compare} found
for the Nosé-Hoover chains and CSVR are within what so few runs leave open,
and the 2.9 of Berendsen coupling lies only just beyond it.
\Cref{sec:cv-observables} returns to those numbers with error bars from
the runs themselves.

\begin{keyidea}
The mean of $n$ independent samples has the spread $\sigma_x/\sqrt n$, and
approaches a Gaussian as $n$ grows whatever the samples' own density,
provided its variance is finite. Its
standard error is $s/\sqrt n$, with $s^2$ the sum of squared deviations
divided by $n - 1$. The interval $\pm1.96$ standard errors holds the true
mean in 95\% of runs when $n$ is large; with few samples it holds it less
often, and Student's factor replaces 1.96.
\end{keyidea}

\notebook{15}{roll dice and watch the mean narrow and the error bars cover}

\begin{selfcheck}
A run of 100 independent samples gives a standard error of 0.05. How many
samples give 0.01?
\end{selfcheck}
